\section{Conventions and static primitives}\label{sec:preliminaries}

\subsection{Representations and fixed-parameter bounds}

Auxiliary graphs may have parallel edges or arcs, always distinguished by
identifiers. Fix a reference orientation on every variable edge. Traversing
an edge backward negates its coefficient. A signed circulation has zero
divergence at every vertex; its length norm for positive lengths $\ell$ is
\[
                       \|\ell f\|_1=\sum_e\ell_e|f_e|.
\]
This norm differs from the \emph{raw length} of a supplied traversal list:
the latter counts every traversal with its multiplicity, before cancellation.
Later a forest edge may map either to an actual graph edge, with its endpoint
correspondence specified, or to a null edge whose two endpoints have the
same image. A null traversal has zero actual projected length, although a
reduction may assign it an additional positive bookkeeping length.

All representation sizes include vertices or slots, individual edges,
terminal incidences, and explicitly stored lists. Ordinary adjacency lists,
heaps, balanced dictionaries, and path-copy persistent maps incur logarithmic
overhead per record operation. No argument permits a free scan of all
coordinates after an implicit path update. Arithmetic and key-length costs
are included separately when they exceed one word.

The matching procedure uses a master bound $S=2^s$, an integer parameter
$a\ge1000$, and a reliability parameter $\sigma$. In the final schedule
$\sigma\le a+1$. A bound of the form
\[
                       S^{O(1/a)+o_a(1)}
\]
has an absolute coefficient in $O(1/a)$; for each fixed $a,\sigma$, the second
exponent tends to zero as $S\to\infty$. Constants and powers of logarithms
may depend on these fixed parameters. All smaller local routines measure
their permitted logarithmic factors against the indicated master bound.
The parameters are not fixed uniformly while taking $a\to\infty$.

An implementation may check a deterministic size or precision prerequisite.
If that check is inadequate, it \emph{suspends} the entire parameter branch
without an answer. In contrast, exhaustion of a capped randomized routine
is a \emph{detected failure}; the calling perfect-matching test returns a
negative answer promptly. The fixed-parameter analysis proves that
deterministic suspensions cannot occur on any path once the input is
sufficiently large for that fixed parameter. It does not require random
attempts always to succeed. These conventions are used in
Theorem~\ref{thm:subcubic} and the schedule in
Section~\ref{sec:uniformization}.

\subsection{Exact integral flow}

\begin{lemma}[Static integral flow]\label{lem:integral-flow}
For an explicitly represented directed network of total size at most $P$,
with integral capacities, demands and costs bounded in magnitude by a fixed
power of $P$, exact integral minimum-cost flow can be computed uniformly and
deterministically in $P^{1+o(1)}$ time. Its finite-precision operations can
be simulated with polynomially logarithmic overhead on logarithmic-bit
words. In particular, integral source--sink maximum flow with nonnegative
upper capacities and zero lower capacities returns an explicit integral
arc vector and a residual-reachability minimum-cut certificate within that
bound.
\end{lemma}

\begin{proof}
Use the deterministic minimum-cost-flow theorem of van den Brand et
al.~\cite[Theorem~1.1]{FlowDeterministic}, with its finite fixed-point model
and deterministic integral rounding~\cite[Section~2 and Lemma~4.1]{FlowDeterministic}.
The stated logarithmic dependence on integral magnitudes and polynomial
logarithmic word lengths is absorbed in $P^{o(1)}$. A fixed-point word is
stored as a charged array of logarithmic-bit words. Longer virtual addresses
can be keys in a dictionary of touched cells; no unused address space is
initialized.

Maximum flow is obtained by adding a bounded sink-to-source return arc and
minimizing its negative flow. Its capacity can be the sum of the original
capacities. Feasibility and absence of a residual source--sink path certify
maximality. Residual reachability supplies the cut.
\end{proof}

Only the nonnegative-upper-capacity maximum-flow specialization is called
algorithmically below. Integral lower bounds are used in feasibility
arguments through the classical cut criteria for circulation
feasibility~\cite{Gale1957,FordFulkerson1958}. The following consequence
records the precision margin we use.

\begin{lemma}[Integral cuts from a small total violation]\label{prelim:small-violation}
Let every lower and upper arc bound be integral, with lower bound no larger
than upper bound. If an exactly conserved real arc vector has total bound
violation less than one, the bounds admit an exactly feasible integral
circulation.
\end{lemma}

\begin{proof}
For a vertex set $X$, feasibility requires the sum of lower bounds on arcs
leaving $X$ to be at most the sum of upper bounds on arcs entering $X$.
Conservation of the real vector implies that any failure of this inequality
is at most its total bound violation. The two bound sums are integers, so
a positive discrepancy would be at least one, a contradiction. The
circulation cut criterion therefore holds. To obtain that criterion, subtract
the lower bounds, introduce a super-source and super-sink for the resulting
vertex imbalances, and apply the integral maximum-flow/minimum-cut theorem.
Saturating the imbalance arcs gives the asserted integral circulation.
\end{proof}

Large capacities relative to a small local graph do not require a different
finite program for each magnitude exponent.

\begin{lemma}[One polynomial-capacity implementation]\label{lem:capacity-scaling}
For maximum flow with nonnegative integral upper capacities of at most $b$
bits, a single fixed polynomial-capacity version of
Lemma~\ref{lem:integral-flow} suffices, using $b$ calls on linear-size
residual graphs, together with the charged capacity arithmetic and identifier
translations.
\end{lemma}

\begin{proof}
Begin with zero capacities. In each binary digit round write the new
capacity vector as $c'=2c+\xi$, where every $\xi_e$ is zero or one.
If $f$ was a maximum flow of value $v$ for $c$, then $2f$ is feasible for
$c'$. Let $A$ be the number of original arcs. An old minimum cut has new
capacity at most $2v+A$, so the largest residual augmentation has integral
value $\Delta\le A$.

Decompose an optimal integral residual augmentation into source--sink paths
and cycles and discard the cycles. Its remaining total path value is
$\Delta$, so every forward and reverse residual arc carries at most
$\Delta\le A$. Truncating every residual capacity at $A$ therefore preserves
the optimum augmentation. Compute it with the fixed polynomial-capacity
routine. Adding forward usage and subtracting reverse usage updates $2f$
to a maximum flow for $c'$. This proves the induction over digits.

Local vertices and arcs can be renumbered, using dictionaries to translate
their caller identifiers on input and output. Integer arrays, comparisons,
and translations are charged at their actual bit or multiword cost.
In the fixed-parameter applications below, $b$ and all key lengths contribute
only the permitted logarithmic powers in the master size.
\end{proof}

When paths from an integral flow are needed, pair incoming and outgoing flow
units and trace source paths, dropping cycles. We charge the total represented
throughput, not merely the number of graph arcs. The routing helper supplies
the explicit throughput bounds at each such use.

\subsection{A clocked fallback}

The final schedule needs one fixed polynomial fallback, solely to guarantee
halting and polynomial physical addressability at every finite input size.
Edmonds' deterministic algorithm supplies the following interface.

\begin{lemma}[Polynomial perfect-matching fallback]\label{lem:fallback}
There is one uniform deterministic algorithm on explicit loopless graphs of
maximum degree three that halts within a fixed polynomial number of finite-word
instructions. It returns a checked perfect matching when one exists and a
negative answer otherwise.
\end{lemma}

\begin{proof}
Run Edmonds' deterministic polynomial-time maximum-cardinality matching
algorithm~\cite[Section~2, pp.~451--452]{EdmondsPaths}. For parallel edge
copies, first deduplicate unordered endpoint pairs while retaining one
original edge identifier for each pair. This preserves perfect matchability
and explicit witness decoding. Return the computed matching if it covers
every vertex, and a negative answer otherwise; check every returned identifier and
vertex incidence before acceptance.

The classical algorithm, the deduplication, and this check have a fixed
polynomial clock without a matchability promise. Their combinatorial records
and polynomially bounded integer indices have polynomial bit cost, so
charged multiword and touched-address implementations retain a fixed
polynomial bound in the stated model. Exact maximum cardinality makes the
returned answer correct in both cases.
\end{proof}

The two witness-preserving reductions used with this fallback are proved
explicitly in Section~\ref{sec:uniformization}.
